\documentclass[12pt]{article}
\usepackage[utf8]{inputenc}
\usepackage[T1]{fontenc}
\usepackage{hyperref}
\title{Khehek Language Resources for K-12 Curriculum Development}
\author{Compiled by Claude Sonnet 5 (claude-sonnet-5) for Andr\'as Kornai}
\date{}
\begin{document}
\maketitle

\section{Introduction}
Khehek (ISO 639-3 tlx; Glottocode kheh1237) is an Oceanic (Austronesian) language of west-central Manus Island, Papua New Guinea, with about 1,600 speakers (a 1991 figure). It has two dialects, Drehet and Levei, and some sources add Bucho. It is not the same as the neighbouring Tulu-Bohuai (rak). Searches found no educational materials.

\section{Primary Educational Resources}
No primers, textbooks or school programmes for Khehek were found.

\section{Linguistic and Reference Materials}
\begin{itemize}
\item Wikipedia, ``Khehek language'': an overview of speakers and dialects. \url{https://en.wikipedia.org/wiki/Khehek_language}
\item Verbix documentation page for Khehek. Its contents were not checked beyond the search listing. \url{https://docs.verbix.com/Languages/Khehek}
\item Rosetta Project (Long Now), entry for tlx. \url{https://rosettapanglossia.longnow.org/index.php?title=Tlx}
\end{itemize}

\section{Religious and Community Resources}
No scripture translations or community organisations were identified. The Global Recordings Network page for tlx may list recordings of a religious nature (see below).

\section{Digital and Audio Resources}
\begin{itemize}
\item Global Recordings Network, Khehek language page. The audio inventory was not verified. \url{https://globalrecordings.net/en/language/tlx}
\end{itemize}

\section{Curriculum Development Context}
Khehek is spoken by a small community, and no school use was found. Any curriculum would need new work: an orthography, a lexicon and recorded texts, made together with the community. Tok Pisin and English are the likely languages of instruction in the region, though the searches did not confirm this.
\end{document}